\documentclass[11pt]{article}
\usepackage{fontspec}
\usepackage[french]{babel}
\usepackage{amsmath}
\usepackage{amssymb}
\usepackage{amsthm}
\usepackage{mathrsfs}
\usepackage{tikz-cd}
\usepackage{enumitem}

\newcommand{\printedpage}[1]{\marginpar{\footnotesize p.~#1}}
\newcommand{\msp}[1]{\marginpar{\footnotesize [#1]}}
\newcommand{\card}{\operatorname{card}}
\newcommand{\Filt}{\mathrm{Filt}}
\newcommand{\Ob}{\operatorname{Ob}}
\newcommand{\circled}[1]{\tikz[baseline=(c.base)]\node[draw,circle,inner sep=1pt](c){#1};}

\begin{document}

\printedpage{44}
% head: A. Grothendieck : Dérivateurs, Ch. XIII, Édition des Universités de Montpellier II et Paris VII

\begin{itemize}
\item[c)] \textit{$I'$ grand devant $\pi_0$.}
\end{itemize}

\texttt{[Le lemme tel qu'il est énoncé est évident pour les mêmes raisons que ci-dessus (même sans l'hypothèse de la note de bas de page). La démonstration qui suit prouve en fait un résultat plus fin :}

\texttt{\textsc{Lemme} 1$'$ : \textit{Soient $\pi_0$, $\pi$ deux cardinaux, avec $\pi_0<\pi$ et $\pi^{\pi_0}=\pi$, $I$ un ensemble ordonné grand devant $\pi_0$, et $I'_0$ une partie de $I$ de cardinal $\le\pi$. Alors il existe une partie $I'$ de $I$ satisfaisant les conditions :}}

\begin{itemize}
\item[\texttt{a)}] $I'\supset I'_0$\texttt{.}
\item[\texttt{b)}] $\card I'\le\pi$\texttt{.}
\item[\texttt{c)}] \texttt{\textit{$I'$ grand devant $\pi_0$. ]}}
\end{itemize}

On va essayer de construire $I'$ par récurrence transfinie, en définissant les $I'_\alpha\subset I$ ainsi :
\[
\begin{aligned}
I'_0 &= I'_0\;, \\
I'_\alpha &= \textstyle\bigcup_{\beta<\alpha} I'_\beta \quad \text{si $\alpha$ est un ordinal limite,} \\
I'_{\alpha+1} &= \left\{\parbox{0.6\textwidth}{déduit de $I'_\alpha$ ainsi : pour toute partie $J$ de $I'_\alpha$ telle que $\card J\le\pi_0$, qui n'est pas majorée dans $I'_\alpha$, on choisit un majorant $i_J$ dans $I$ (\texttt{[qui]} existe, car $I$ \texttt{[est]} grand devant $\pi$, donc devant $\pi_0$) et $I'_{\alpha+1}$ s'obtient en adjoignant ces $i_J$ à $I'_\alpha$.}\right.
\end{aligned}
\]

\msp{page 81}
Soit alors $\alpha_0$ le premier ordinal qui soit de cardinal $>\pi_0$, et considérons
\[
I' \;=\; I'_{\alpha_0} \;=\; \bigcup\limits_{\substack{\alpha<\alpha_0\\ \textit{i.e.}\ \card\alpha\le\pi_0}} I'_\alpha .
\]
Alors $I'$ contient $I'_0$, il est grand devant $\pi_0$, mais quel est son cardinal ? Prouvons par récurrence transfinie que
\[
\card I'_\alpha \;\le\pi \qquad\quad \text{si } \alpha<\alpha_0,\ \ \textit{i.e.}\ \ \card\alpha\le\pi_0 .
\]
C'est O.K. pour $\alpha=0$ par hypothèse sur $I'_0$. Si c'est vrai pour $I'_\alpha$, est-ce vrai pour $I'_{\alpha+1}$ ? On adjoint à $I'_\alpha$ un ensemble d'éléments $i_J$ qui est de cardinal $\le\card\mathfrak{P}_{\le\pi_0}(I'_\alpha)$, donc de cardinal $\le\card(I'_\alpha)^{\pi_0}\le\pi^{\pi_0}$, et si on suppose
\[
\pi^{\pi_0} \;=\; \pi,
\]
on gagne. Supposons enfin que $\alpha$ soit un ordinal limite (de cardinal $\le\pi_0$). Donc on a
\[
I'_\alpha \;=\; \bigcup_{\beta<\alpha} I'_\beta ,
\]
donc
\printedpage{45}
% head: A. Grothendieck : Dérivateurs, Ch. XIII, Édition des Universités de Montpellier II et Paris VII
\msp{page 82}
\[
\card I'_\alpha \;\le\; \underbrace{(\card\alpha)\pi}_{\le\pi\pi=\pi^2=\pi} \qquad\quad \text{puisque les } \card I'_\beta\le\pi,
\]
on gagne encore.

On trouve enfin, par le même argument,
\[
\card I'_{\alpha_0} \;\le\; \underbrace{(\card\alpha_0)\pi\le\pi^2}_{\substack{\text{si }\card\alpha_0\le\pi,\\ \textit{i.e.}\text{ si }\pi_0<\pi}} = \pi,
\]
donc on gagne encore.

Ainsi, on a trouvé un deuxième système inductif
\[
(X_{\tilde\imath})_{\tilde\imath\in\tilde I}
\]
du même type que le système initial, ayant comme limite inductive $X$ (qu'il ait même limite inductive est immédiat). Mais qu'a-t-on gagné ? Ceci : dans $\tilde I$ toute partie de cardinal $\le\pi$ et grande devant $\pi_0$ est non seulement \textit{majorée} dans $\tilde I$, mais même y admet une borne supérieure. Cela résulte du % margin: double vertical rule beside the four lines "inductive est immédiat" ... "supérieure. Cela résulte du"

\msp{page 83}
\textbf{Lemme 2 :} \textit{Soit $I$ grand devant $\pi>\pi_0$, avec $\pi^{\pi_0}=\pi$} \texttt{[il suffit $\pi\ge\pi_0$]} \textit{(cf. lemme 1). Soit $(I'_\alpha)_{\alpha\in A}$ un ensemble de parties de $I$ qui sont dans $\tilde I$,} i.e.

\begin{itemize}
\item[a)] \textit{de cardinal $\le\pi$, et}
\item[b)] \textit{grandes devant $\pi_0$,}
\end{itemize}

\textit{et supposons que} $\card A\le\pi$ \textit{et que $A$, ordonné par inclusion, soit grand devant $\pi_0$. Alors $I'=\bigcup_{\alpha\in A} I'_\alpha$ est encore dans $\tilde I$,} i.e. \textit{$I'$ est de cardinal $\le\pi$ (car de cardinal $\le\pi^2=\pi$), et de plus, grande devant $\pi_0$.} (Vérification immédiate.)

Ainsi, on a prouvé ceci :

% margin: single vertical rule beside Proposition 1, from here to the end of the page
\textbf{Proposition 1.} \textit{Soit $E$} \texttt{[une]} \textit{grosse catégorie, avec filtration cardinale} $(\Filt^\pi)_{\pi\ge\pi_0}$. \textit{Soit $\pi$ un cardinal avec}
\[
\pi>\pi_0, \qquad\quad \pi^{\pi_0}=\pi \qquad\quad \textit{(p. ex. $\pi=2^c$, avec $c\ge\pi_0$)} \quad \texttt{[}\pi\ge\pi_0\ \texttt{suffit].}
\]
\textit{Alors pour tout $X$ dans $E$, on peut trouver un système inductif $(X_i)_{i\in I}$, de limite $X$, avec}

\begin{itemize}
\item[a)] \[ \left\{ \begin{aligned} & X_i\in\Filt^\pi \ \ \textit{pour tout $i\in I$,} \\ & \ I \ \textit{grand devant $\pi$,} \end{aligned} \right. \]
\end{itemize}

\textit{satisfaisant aux conditions supplémentaires suivantes :}

\begin{itemize}
\item[b)] \textit{Toute partie $J$ de $I$ de cardinal $\le\pi$, et grande devant $\pi_0$, admet une borne supérieure $i_J$ dans $I$.}
\item[c)] \textit{On a, sous ces conditions}
\end{itemize}

\printedpage{46}
% head: A. Grothendieck : Dérivateurs, Ch. XIII, Édition des Universités de Montpellier II et Paris VII
\msp{page 84}
\[
X_{i_J} \xleftarrow{\ \sim\ } \varinjlim_{i\in J} X_i .
\]
\textit{De plus, si $X\in\Filt^{\pi'}$ avec $\pi'>\pi$, on peut prendre $I$ de cardinal $\le\pi'^{\pi}$ (donc de cardinal $\le\pi'$ si $\pi'^{\pi}=\pi'$, p. ex. si $\pi'=2^{c'}$, avec $c'\ge\pi$).}

\textbf{Définition.} On appelle un système inductif $(X_i)_{i\in I}$, satisfaisant les conditions \textit{a}), \textit{b}), \textit{c}) ci-dessus, un \textit{système inductif bien adapté} (à la filtration cardinale donnée, pour les cardinaux $\pi_0$, $\pi$ donnés), et la représentation correspondante
\[
X \xleftarrow{\ \sim\ } \varinjlim X_i
\]
une \textit{représentation bien adaptée} de $X$ en termes d'objets de $\Filt^\pi$.

\textbf{NB} Quand $E$ est stable par petites $\varinjlim$ filtrantes, la situation se simplifie considérablement. On peut alors prendre
\[
\tilde I \;=\; \{I'\subset I \mid \card I'\le\pi,\ I' \text{ filtrant}\}
\]
(sans s'encombrer de la condition $I'$ grand devant $\pi_0$) et définir encore le système inductif $(X_{I'})_{I'\in\tilde I}$ par
\[
X_{I'} \;=\; \varinjlim_{i\in I'} X_i \;.
\]
On a encore
\msp{page 85}
\[
\card\tilde I \;\le\; (\card I)^\pi \;\le\; \pi'^{\pi} \qquad\quad \text{si }\ \card I\le\pi'^{\pi}
\]
(notamment, si $X\in\Filt^{\pi'}$, et qu'on prenne $I$ comme indiqué dans la condition $c$ des filtrations cardinales). De plus, cette fois, le fait que $\tilde I$ soit grand devant $\pi$ est trivial, on a même mieux, à part que $I$ soit filtrante :

\begin{itemize}
\item[$(*)$] \textit{Toute} partie $(I'_\alpha)_{\alpha\in A}$ de $\tilde I$ filtrante et de cardinal $\le\pi$ admet une borne supérieure $\tilde\imath_A$,
\end{itemize}

savoir la réunion (qui est bien filtrante). De plus, le système inductif $(X_{\tilde\imath})_{\tilde\imath\in\tilde I}$ a la propriété plus forte que
\[
X_{\tilde\imath_A} \xleftarrow{\ \sim\ } \varprojlim_{\tilde\imath\in A} X_{\tilde\imath} \;. % [sic] \varprojlim (arrow under lim pointing left) -> \varinjlim, inductive limit of the inductive system X_ĩ over the filtrant part A
\]
Ainsi on a la

\textbf{Proposition 1 bis.} \textit{Supposons $E$ stable par petites $\varinjlim$ filtrantes, et} $\Filt^\pi$ \textit{stable par celles-ci, où $\pi\ge\pi_0$ est donné. Alors pour tout $X$ dans $E$, on peut trouver un système inductif $(X_i)_{i\in I}$ de limite inductive $X$, avec}

\begin{itemize}
\item[a)] $X_i\in\Filt^\pi$ \textit{pour tout $i\in I$, $I$ ensemble ordonné grand devant $\pi$.}
\item[b)] \textit{Toute partie $J$ de $I$ filtrante de cardinal $\le\pi$ admet une borne supérieure $i_J$ dans $I$.}
\item[c)] \textit{Sous ces conditions}
\end{itemize}

\printedpage{47}
% head: A. Grothendieck : Dérivateurs, Ch. XIII, Édition des Universités de Montpellier II et Paris VII
\msp{page 86}
\[
X_{i_J} \;=\; \varinjlim_{i\in J} X_i \;.
\]
De tels systèmes inductifs seront dits \textit{“parfaitement adaptés”} (à la filtration cardinale donnée, pour le cardinal $\pi$ donné – $\pi_0$ n'y figure plus, si ce n'est comme début de la filtration cardinale sans plus).

Comme pratiquement, on aura à manier surtout des $E$ stables par petites $\varinjlim$ filtrantes, et des filtrations cardinales stables en conséquence (renforcement de la condition \textit{b}) des filtrations cardinales), c'est le cas, et la notion de \textit{représentation parfaitement adaptée} de $X$ comme $\varinjlim\limits_{I} X_i$ ($X_i\in\Filt^\pi$), que je mettrai en avant. Le cas “général”, où on suppose seulement \texttt{[que]} $E$ satisfait $L_{\pi_0}$, est plus alambiqué et se traite moins joliment (\footnote[41]{On parlera de filtration cardinale \textit{stricte,} si les $\underline{\Filt}^\pi$ sont stables par $\varinjlim$ filtrantes de cardinal $\le\pi$ (et que ces limites filtrantes petites existent dans $E$).}).

\textbf{Définition 2.} Soit $E$ munie d'une filtration cardinale stricte $(\underline{\Filt}^\pi)_{\pi\ge\pi_0}$. Soit $M$ une partie de $\Ob E$, $\underline{M}$ la sous-catégorie pleine de $E$ qu'elle définit, et soit $\pi$ un cardinal $\ge\pi_0$. On dit que $M$ est \textit{$\pi$-accessible,} où que $\underline{M}$ est \textit{$\pi$-accessible}, % [sic] où que -> ou que, disjunction "ou" not the relative "où"
\msp{page 87}
si les conditions suivantes sont satisfaites. Pour tout système inductif
\[
(X_i)_{i\in I}
\]
parfaitement adapté à la filtration cardinale pour $\pi$ (donc $I$ ensemble ordonné filtrant tel que toute partie filtrante $J$ de cardinal $\le\pi$ ait une borne supérieure $i_J$, et $X_{i_J}\xleftarrow{\ \sim\ }\varinjlim\limits_{j\in J} X_j$) ; soit $I_M$ l'ensemble des $i\in I$ tels que $X_i\in M$. Ceci posé, on veut ceci : pour que $X$ ($=\varinjlim X_i$) soit dans $M$, il faut et il suffit que $I_M$ soit cofinal dans $I$ (de sorte que $X$ apparaît comme $\varinjlim$ filtrante, avec ensemble d'indices $I_M$ grand devant $\pi$, d'éléments de $M\cap\underline{\Filt}^\pi\overset{\text{déf}}{=}M^\pi$).

\msp{page 88}
J'ai envie de reprendre la définition de la $\pi$-accessibilité de $M$, pour une filtration cardinale donnée, à coups de conditions équivalentes. Pour tout système inductif $(X_i)_{i\in I}$ dans $M$ \texttt{[plutôt dans $E$]}, je garde la notation
\[
I_M \;=\; \{i\in I \mid X_i\in M\}.
\]
Je pose, pour tout cardinal $\pi\ge\pi_0$,
\[
M^\pi \;=\; M\cap\Filt^\pi .
\]

\begin{itemize}
\item[\circled{a$_\pi$}] $M$ est l'ensemble des objets de la forme
\[
\varinjlim_{i\in I} X_i \;,
\]
où les $X_i$ sont dans $M^\pi$, et $I$ est grand devant $\pi$.
\end{itemize}

\printedpage{48}
% head: A. Grothendieck : Dérivateurs, Ch. XIII, Édition des Universités de Montpellier II et Paris VII

\begin{itemize}
\item[\circled{b$_\pi$}] Soit $(X_i)_{i\in I}$ un système inductif dans $M$ \texttt{[plutôt dans $E$]} parfaitement adapté à la filtration cardinale pour $\pi$, \textit{i.e.} $I$ filtrante et toute partie filtrante $J$ de $I$ de cardinal $\le\pi$ admet une borne supérieure $i_J$, et $X_{i_J}\xleftarrow{\ \sim\ }\varinjlim\limits_{J} X_i$, et enfin $X_i\in\Filt^\pi$ pour tout $i\in I$. Soit $X=\varinjlim X_i$. Pour que l'on ait $X\in M$, il faut et il suffit que $I_M$ soit cofinal dans $I$.
\item[\circled{c$_\pi$}] $\underline{M}^\pi=\underline{M}\cap\underline{\Filt}^\pi$ est une sous-catégorie de $\underline{M}$ génératrice par épimorphismes stricts, $\underline{M}$ est stable par $\varinjlim$ grandes devant $\pi$, enfin pour tout $X\in M$, $\underline{M}^\pi/X$ est une catégorie filtrante.
\end{itemize}

\msp{page 89}
Je présume que ces conditions sont équivalentes.

$b_\pi\Longrightarrow a_\pi$. On suppose \textit{b}) vérifié, on veut prouver \textit{a}). D'abord que tout $X\in M$ est de la forme dite $\varinjlim\limits_{I} X_i$, avec les $X_i\in M^\pi$ et $I$ grand devant $\pi$. Par \texttt{[la]} proposition 1, on sait qu'on a une présentation parfaite adaptée pour $\pi$, $X=\varinjlim\limits_{i\in I} X_i$. Par \textit{b}) “il faut”, la partie $I_M$ de $I$ est cofinale dans $I$, donc comme $I$ est grand devant $\pi$, $I_M$ aussi, et on aura (posant $J=I_M$)
\[
X \;=\; \varinjlim_{i\in J} X_i, \qquad\quad X_i\in M^\pi,\ \ J \text{ grand devant } \pi,
\]
O.K. Inversement, prouvons que si $X$ se met sous la forme dite dans \textit{a}), il est dans $M$. Cela va résulter du

\textbf{Lemme 1 :} \textit{Si la condition} b) (\textit{ou} a)) \textit{est satisfaite, $M$ est stable par limites inductives grandes devant $\pi$.}

Il faut prouver que si
\msp{page 90}
\[
X \;=\; \varinjlim_{i\in I} X_i \;,
\]
avec $I$ grand devant $\pi$, les $X_i$ dans $M$, alors $X\in M$. La condition \textit{b}) “il suffit” l'assure quand on suppose de plus les $X_i$ dans $\Filt^\pi$ (donc dans $M^\pi$), de sorte que l'on aura $I_M=I$, donc $I_M$ cofinal, donc $X\in M$ (\footnote[42]{Pas tout à fait – on y suppose plus, sur le système inductif. On ne peut se dispenser, même dans le cas $X_i\in M^\pi$ $\forall\, i\in I$, faire le raisonnement des pages suivantes.}). Si on ne suppose pas les $X_i$ dans $\Filt^\pi$, il faut travailler un peu plus.

Passer tout de suite à la partie de la démonstration ci-dessous, page 96, qui prouve \textit{b}) $\Longrightarrow$ \textit{a}). Le reste de la démonstration va prouver le

\textbf{Lemme 2.} \textit{Si tout élément de Filt$^\pi$ est $\pi$-accessible, alors} b$_\pi$) \textit{(et même} a$_\pi$\textit{) implique que $M$ est stable par $\varinjlim$ grandes devant $\pi$.}

\end{document}
